\scanpage{150}
\mspage{74}
Considérons le système $\frac12$ simplicial des
\[
(245)\quad
\diagram{\begin{tikzcd}[ampersand replacement=\&,column sep=small]
\& \mathcal V^{*}_{d_0\ldots d_r}\mathrel{\overset{\mathrm{déf}}=}\mathcal V^{*}_{D_{d_0,\ldots,d_r},X}\arrow[d,hook] \& \text{où }X=D_0\\
D_{d_0,\ldots,d_r}\arrow[r] \& \mathcal V_{d_0,\ldots,d_r}=\mathcal V_{D_{d_0,\ldots,d_r},X} \&
\end{tikzcd}}
\]
\margin{On se \unsure{suppose} $d_0\geq d_1\geq\cdots\geq d_r\unsure{\geq|>}0$.}
les $\mathcal V^{*}_{d_0,\ldots,d_r}$ s'envoient tous \deleted{\unsure{les uns}} \inserted{dans $X^{*}=D_0^{*}=X\setminus\Theta$} et sauf $\mathcal V^{*}_0=X^{*}$, ils \deleted{via $\mathcal V_d^{*}\lhook\joinrel\longrightarrow D$} s'envoient dans $\mathcal V^{*}_{\Theta,X}$.

\uline{Th de recollement}\quad Considérons le \deleted{\unsure{fibré en}} topos \inserted{$\frac12$ simplicial au dessus de $\mathcal V^{*}_{\Theta,X}$, \deleted{\unsure{dont les composantes sont}}} (formé de \unsure{multiplicités} analytiques) formé des $\mathcal V^{*}_{d_0,\ldots,d_r}$ ($d_0\geq d_1\ldots\geq d_r>0$), \deleted{\unsure{Alors on}} soit $\mathcal V^{*}_{\unsure{\bullet}}$, on a un morphisme de la limite inductive de ces topos dans $\mathcal V^{*}_{\Theta,X}$ \inserted{$=X^{*\infty}=D_0^{*\infty}$}, et c'est là une équivalence de topos (a fortiori, \unsure{c'est} une équivalence d'homotopie, et en particulier ça induit un isom. sur les $\pi_0$, et les $\pi_1$ relatifs à des comp. connexes qui se correspondent).
\margin{149}
